\documentclass[a4paper,10pt,twocolumn]{ltjarticle}
\usepackage{kaitoushu}

\begin{document}

\problemrange{4章：指数関数と対数関数 — 練習問題1・B (p.111)}



\mondai{1}{$x+x^{-1}=18$ を満たす正の実数 $x$ に対して，次の式の値を求めよ．}
\kotae{
\textbf{(1)} $x^{1/2}+x^{-1/2}$

$t=x^{1/2}+x^{-1/2}>0$ とおくと
\[
t^{2}=x+2+x^{-1}=18+2=20
\]

$\therefore t=2\sqrt{5}$

\textbf{(2)} $x^{1/3}+x^{-1/3}$

$s=x^{1/3}+x^{-1/3}\;(>0)$ とおき3乗する：
\[
s^{3}=x+3(x^{1/3}\cdot x^{-1/3})s+x^{-1}=18+3s
\]

よって $s^{3}-3s-18=0$．$s=3$ が解（$27-9-18=0$）で，
\[
(s-3)(s^{2}+3s+6)=0
\]

$s^{2}+3s+6$ は判別式 $9-24<0$ より実数解なし．

$\therefore s=3$
}

\mondai{2}{次の数を小さいものから順に並べよ．}
\kotae{
\textbf{(1)} $\dfrac{1}{\sqrt[3]{16}},\;\dfrac{\sqrt[8]{2^{-1}}}{2},\;8^{-2/5}$

底を $2$ に統一：
\[
\dfrac{1}{\sqrt[3]{16}}=2^{-4/3},\;\dfrac{\sqrt[8]{2^{-1}}}{2}=2^{-9/8},\;8^{-2/5}=2^{-6/5}
\]

指数の大小：$-\dfrac{4}{3}\approx-1.333,\;-\dfrac{6}{5}=-1.2,\;-\dfrac{9}{8}=-1.125$．

$-\dfrac{4}{3}<-\dfrac{6}{5}<-\dfrac{9}{8}$．底 $2>1$ は単調増加．

$\therefore \dfrac{1}{\sqrt[3]{16}}<8^{-2/5}<\dfrac{\sqrt[8]{2^{-1}}}{2}$

\textbf{(2)} $\sqrt[4]{3^{2}},\;\sqrt[4]{10},\;(\sqrt[4]{2})^{3}$

すべて4乗根の形に：
\[
\sqrt[4]{9},\;\sqrt[4]{10},\;\sqrt[4]{8}
\]

$8<9<10$ より

$\therefore (\sqrt[4]{2})^{3}<\sqrt[4]{3^{2}}<\sqrt[4]{10}$
}

\mondai{3}{$x>0$ とする．$3^{x}+3^{-x}=3$ のとき，$3^{x}-3^{-x}$ の値を求めよ．}
\kotae{
\[
(3^{x}-3^{-x})^{2}=(3^{x}+3^{-x})^{2}-4=9-4=5
\]

$x>0$ より $3^{x}>3^{-x}$ なので $3^{x}-3^{-x}>0$．

$\therefore 3^{x}-3^{-x}=\sqrt{5}$
}

\mondai{4}{次の不等式および方程式を解け．}
\kotae{
\textbf{(1)} $4^{x}-3\cdot2^{x}+2<0$

$t=2^{x}>0$ とおくと $t^{2}-3t+2<0$，$(t-1)(t-2)<0$．
$1<t<2$，すなわち $2^{0}<2^{x}<2^{1}$．

$\therefore 0<x<1$

\textbf{(2)} $2\cdot4^{2x+1}-6\cdot4^{x}+1=0$

$4^{2x+1}=4(4^{x})^{2}$．$t=4^{x}>0$ とおくと
\[
8t^{2}-6t+1=0 \Rightarrow (4t-1)(2t-1)=0
\]

$t=\dfrac{1}{4}$：$4^{x}=4^{-1} \Rightarrow x=-1$．
$t=\dfrac{1}{2}$：$4^{x}=4^{-1/2} \Rightarrow x=-\dfrac{1}{2}$．

$\therefore x=-1,\;x=-\dfrac{1}{2}$

\textbf{(3)} $4^{x}-7\cdot2^{x}\leqq 8$

$t=2^{x}>0$ とおくと $t^{2}-7t-8\leqq0$，$(t-8)(t+1)\leqq0$．
$t>0$ より $0<t\leqq8$，すなわち $2^{x}\leqq2^{3}$．

$\therefore x\leqq3$

\textbf{(4)} $\begin{cases} 2^{x+1}+2^{1-y}=9 \\ 2^{x-y+1}=4 \end{cases}$

第2式：$x-y+1=2 \Rightarrow x=y+1$．

第1式に代入：$2^{y+2}+2^{1-y}=9$．$t=2^{y}>0$ とおくと
\[
4t+\dfrac{2}{t}=9 \Rightarrow 4t^{2}-9t+2=0
\]
\[
(4t-1)(t-2)=0 \Rightarrow t=\dfrac{1}{4},\;2
\]

$t=2$：$y=1,\;x=2$．
$t=\dfrac{1}{4}$：$y=-2,\;x=-1$．このとき第1式
$2^{x+1}+2^{1-y}=2^{0}+2^{3}=1+8=9$ で成立．

$\therefore (x,y)=(2,1),\;(-1,-2)$
}

\mondai{5}{次の不等式を解け．ただし，$a>0,\;a\neq1$ とする．
\[
\dfrac{1}{a^{3x+2}}>\sqrt[3]{a^{2}}
\]}
\kotae{
$a^{-(3x+2)}>a^{2/3}$

\textbf{(i) $a>1$ のとき}：単調増加なので
$-(3x+2)>\dfrac{2}{3} \Rightarrow -3x>\dfrac{8}{3} \Rightarrow x<-\dfrac{8}{9}$

\textbf{(ii) $0<a<1$ のとき}：単調減少なので不等号反転
$-(3x+2)<\dfrac{2}{3} \Rightarrow x>-\dfrac{8}{9}$
}

\mondai{6}{$a>0,\;b>0$ のとき，次の式を簡単にせよ．}
\kotae{
\textbf{(1)} $(\sqrt[4]{a}+\sqrt[4]{b})(\sqrt{a}-\sqrt[4]{ab}+\sqrt{b})$

$A=a^{1/4},\;B=b^{1/4}$ とおくと
\[
(A+B)(A^{2}-AB+B^{2})=A^{3}+B^{3}
\]

$\therefore a^{3/4}+b^{3/4}$

\textbf{(2)} $\dfrac{\sqrt{ab^{3}}}{\sqrt[3]{a^{2}b^{4}}}\times\sqrt[6]{3\times\left(\dfrac{a}{3b}\right)^{-5}}$

第1因子：
\[
\dfrac{a^{1/2}b^{3/2}}{a^{2/3}b^{4/3}}=a^{-1/6}b^{1/6}
\]

第2因子：$\left(\dfrac{a}{3b}\right)^{-5}=\dfrac{3^{5}b^{5}}{a^{5}}$．
\[
\sqrt[6]{3\cdot\dfrac{3^{5}b^{5}}{a^{5}}}=\sqrt[6]{\dfrac{3^{6}b^{5}}{a^{5}}}=3\cdot a^{-5/6}b^{5/6}
\]

両者の積：
\[
a^{-1/6}b^{1/6}\cdot3a^{-5/6}b^{5/6}=3a^{-1}b=\dfrac{3b}{a}
\]
}

\mondai{7}{$a>0$ のとき，次の不等式が成り立つことを証明せよ．
\[
\dfrac{\sqrt{a^{x}}+\sqrt{a^{y}}}{2}\geqq a^{(x+y)/4}
\]}
\kotae{
$\sqrt{a^{x}}>0,\;\sqrt{a^{y}}>0$ なので相加平均$\geqq$相乗平均：
\[
\dfrac{\sqrt{a^{x}}+\sqrt{a^{y}}}{2}\geqq\sqrt{\sqrt{a^{x}}\cdot\sqrt{a^{y}}}
\]
\[
=\sqrt{a^{(x+y)/2}}=a^{(x+y)/4}
\]

等号成立は $\sqrt{a^{x}}=\sqrt{a^{y}}$，すなわち $x=y$ のとき．\hfill$\blacksquare$
}

\mondai{8}{$2x+3y=1$ のとき，$4^{x}+8^{y}$ の最小値を求めよ．}
\kotae{
$4^{x}=2^{2x}>0,\;8^{y}=2^{3y}>0$ なので相加$\geqq$相乗：
\[
4^{x}+8^{y}=2^{2x}+2^{3y}\geqq2\sqrt{2^{2x}\cdot2^{3y}}
\]
\[
=2\cdot2^{(2x+3y)/2}=2\cdot2^{1/2}=2\sqrt{2}
\]

等号成立は $2^{2x}=2^{3y}$，すなわち $2x=3y$．
$2x+3y=1$ と合わせて $2x=3y=\dfrac{1}{2}$，
$x=\dfrac{1}{4},\;y=\dfrac{1}{6}$．

$\therefore$ 最小値 $2\sqrt{2}\quad\left(x=\dfrac{1}{4},\;y=\dfrac{1}{6}\right)$
}

\end{document}
